\begin{proof}[Proof of \cref{obj:lem:observable-odds}]
Fix \(P\in\mathcal M(\alpha,\beta)\).

\begin{enumerate}
\item Define, for \(a,y\in\{0,1\}\) and \(x\in[0,1]\),
\[
p_{ay}(P,x)=P(A=a,Y=y\mid X=x),\qquad
\mu_a(P,x)=\frac{p_{a1}(P,x)}{p_{a0}(P,x)+p_{a1}(P,x)},
\]
and
\[
m_{\mathrm{obs}}(P,x)=E_P[Y\mid X=x].
\]
The strictly positive, normalized conditional cells give
\[
0<e(P,x)<1,\qquad 0<\mu_a(P,x)<1,
\]
and
\[
\begin{aligned}
e(P,x)&=p_{10}(P,x)+p_{11}(P,x),\\
m_{\mathrm{obs}}(P,x)&=p_{01}(P,x)+p_{11}(P,x),\\
p_{11}(P,x)&=e(P,x)\mu_1(P,x),&
p_{10}(P,x)&=e(P,x)\{1-\mu_1(P,x)\},\\
p_{01}(P,x)&=\{1-e(P,x)\}\mu_0(P,x),&
p_{00}(P,x)&=\{1-e(P,x)\}\{1-\mu_0(P,x)\}.
\end{aligned}
\]
Indeed, each arm risk is the positive success cell divided by its arm total; normalization identifies the control-arm total with \(1-e(P,x)\). Substitution gives the four factorizations.

Normalization also gives
\[
p_{00}(P,x)=1-p_{11}(P,x)-p_{10}(P,x)-p_{01}(P,x).
\]
Expanding the product of the two margins therefore yields
\[
p_{11}(P,x)-e(P,x)m_{\mathrm{obs}}(P,x)
=p_{11}(P,x)p_{00}(P,x)-p_{10}(P,x)p_{01}(P,x).
\]

For the logistic map, write
\[
\ell(\vartheta)=\frac{1}{1+\exp(-\vartheta)}
\qquad(\vartheta\in\mathbb R).
\]
Its odds and log-odds satisfy
\[
\frac{\ell(\vartheta)}{1-\ell(\vartheta)}=\exp(\vartheta),
\qquad
\operatorname{logit}\ell(\vartheta)=\vartheta.
\]
Since \(\nu(P,x)=\operatorname{logit}\mu_0(P,x)\), the homogeneous relation supplied by \cref{obj:ass:homogeneous-logit} gives
\[
\operatorname{logit}\mu_1(P,x)
-\operatorname{logit}\mu_0(P,x)=\theta(P).
\]
Exponentiating this identity at the interior arm risks gives
\[
\exp\{\theta(P)\}
=
\frac{\mu_1(P,x)/\{1-\mu_1(P,x)\}}
     {\mu_0(P,x)/\{1-\mu_0(P,x)\}}.
\]
Using the four cell factorizations and clearing the positive denominators,
\[
p_{11}(P,x)p_{00}(P,x)
=\exp\{\theta(P)\}\,p_{10}(P,x)p_{01}(P,x).
\]

The cell disintegration and the uniform design in \cref{obj:ass:uniform-design} imply
\[
E_P[AY]=\int_{[0,1]}p_{11}(P,x)\,\lambda(dx).
\]
Here the binary product \(AY\) selects exactly the cell \(A=Y=1\). Continuity on the compact covariate interval ensures integrability of the cell functions and their products. Consequently,
\[
\begin{aligned}
C(P)
&=\int_{[0,1]}
 \{p_{11}(P,x)-e(P,x)m_{\mathrm{obs}}(P,x)\}\,\lambda(dx)\\
&=\int_{[0,1]}
 \{p_{11}(P,x)p_{00}(P,x)-p_{10}(P,x)p_{01}(P,x)\}\,\lambda(dx)\\
&=\{\exp(\theta(P))-1\}
  \int_{[0,1]}p_{10}(P,x)p_{01}(P,x)\,\lambda(dx)\\
&=\{\exp(\theta(P))-1\}S(P).
\end{aligned}
\]
% lean: CausalSmith.Stat.LogoddsLowsmoothFrontier.covarianceCoordinate_eq_odds

\item We next establish the numerical denominator floor. The numerical exponential bound gives
\[
\exp(1)<2.7182818286<\frac{2719}{1000}<\frac{11}{4}.
\]
% lean: Real.exp_one_lt_d9
Moreover,
\[
\exp(3/2)^2=\exp(1)^3
<\left(\frac{2719}{1000}\right)^3
<\frac{81}{4}.
\]
Positivity of the exponential implies \(\exp(3/2)<9/2\). Hence
\[
\ell(-1)=\frac{1}{1+\exp(1)}\ge\frac4{15},
\qquad
\ell(-3/2)=\frac{1}{1+\exp(3/2)}\ge\frac2{11}.
\]

The bounds in
\cref{obj:ass:propensity-envelope,obj:ass:prognosis-envelope,obj:ass:effect-envelope}
supply, at every \(x\),
\[
-1\le g(P,x)\le1,\qquad
-1\le\nu(P,x)\le1,\qquad
-\frac12\le\theta(P)\le\frac12.
\]
Inverting the native logits at the interior probabilities gives
\[
e(P,x)=\ell\{g(P,x)\},\qquad
\mu_0(P,x)=\ell\{\nu(P,x)\}.
\]
The logistic map is increasing and satisfies
\[
1-\ell(\vartheta)=\ell(-\vartheta)
\qquad(\vartheta\in\mathbb R).
\]
Together with the homogeneous relation, these facts yield
\[
\begin{aligned}
e(P,x)&\ge\ell(-1)\ge\frac4{15},\\
1-e(P,x)&\ge\ell(-1)\ge\frac4{15},\\
\mu_0(P,x)&\ge\ell(-1)\ge\frac4{15}\ge\frac2{11},\\
1-\mu_1(P,x)
&=\ell\{-\nu(P,x)-\theta(P)\}
 \ge\ell(-3/2)\ge\frac2{11}.
\end{aligned}
\]
Multiplying these nonnegative lower bounds,
\[
\begin{aligned}
p_{10}(P,x)p_{01}(P,x)
&=e(P,x)\{1-\mu_1(P,x)\}\{1-e(P,x)\}\mu_0(P,x)\\
&\ge\frac4{15}\frac2{11}\frac4{15}\frac2{11}
=\frac{64}{27225}
>\frac1{512}=s_0.
\end{aligned}
\]
Since \(\lambda\) is a probability law, integration gives
\[
S(P)\ge\int_{[0,1]}s_0\,\lambda(dx)=s_0.
\]
% lean: CausalSmith.Stat.LogoddsLowsmoothFrontier.denominatorCoordinate_lower

\item In particular, \(S(P)\ge s_0>0\), so the coordinate identity permits division:
\[
1+\frac{C(P)}{S(P)}
=1+\exp\{\theta(P)\}-1
=\exp\{\theta(P)\}>0.
\]
Applying the logarithm proves
\[
\theta(P)=\log\{1+C(P)/S(P)\}.
\]
% lean: CausalSmith.Stat.LogoddsLowsmoothFrontier.effect_eq_log_of_coordinates

\item Fix an admissible kernel \(\Gamma\). Define its full conditional law by
\[
\mathcal K_{P,\Gamma}(x)
=\Gamma(P,x)\otimes\operatorname{Bernoulli}\{e(P,x)\},
\qquad x\in[0,1],
\]
and define the reassociation map by
\[
\begin{aligned}
\mathsf r_{\mathrm{assoc}}:\;
[0,1]\times(\{0,1\}^2\times\{0,1\})
&\longrightarrow([0,1]\times\{0,1\}^2)\times\{0,1\},\\
\mathsf r_{\mathrm{assoc}}\bigl(x,((y^0,y^1),a)\bigr)
&=\bigl((x,(y^0,y^1)),a\bigr).
\end{aligned}
\]
For fixed \(x\), the treatment kernel assigns
\(\operatorname{Bernoulli}\{e(P,x)\}\) to every potential-outcome pair. Its composition with \(\Gamma(P,x)\) is therefore precisely \(\mathcal K_{P,\Gamma}(x)\). This also makes \(x\mapsto\mathcal K_{P,\Gamma}(x)\) a measurable probability kernel.

Associativity of successive kernel composition, applied to the construction in \cref{obj:def:causal-extension}, now gives
\[
\widetilde P_{P,\Gamma}
=(\mathsf r_{\mathrm{assoc}})_\#
 \bigl\{\lambda(dx)\,
 \mathcal K_{P,\Gamma}(x)(d((y^0,y^1),a))\bigr\}.
\]
Here the measure inside braces is the joint law obtained by drawing \(x\) from \(\lambda\) and then \(((y^0,y^1),a)\) from \(\mathcal K_{P,\Gamma}(x)\), and the subscript \(\#\) denotes pushforward. This proves the asserted full conditional disintegration.
% lean: CausalSmith.Stat.LogoddsLowsmoothFrontier.observable_odds

\item Define the observation map by
\[
\begin{aligned}
\mathsf o_{\mathrm{obs}}:\;
([0,1]\times\{0,1\}^2)\times\{0,1\}
&\longrightarrow[0,1]\times\{0,1\}^2,\\
\mathsf o_{\mathrm{obs}}\bigl((x,(y^0,y^1)),a\bigr)
&=\bigl(x,a,(1-a)y^0+ay^1\bigr).
\end{aligned}
\]
Let
\[
\varphi:[0,1]\times\{0,1\}^2\longrightarrow[0,\infty]
\]
be any nonnegative measurable test function. At each fixed \(x\), expand the treatment Bernoulli law into its two atoms. The first margin of \(\Gamma(P,x)\) handles \(a=0\), and the second handles \(a=1\). Thus
\[
\begin{aligned}
&\int_{\{0,1\}^2}\int_{\{0,1\}}
 \varphi\bigl(x,a,(1-a)y^0+ay^1\bigr)\,
 \operatorname{Bernoulli}\{e(P,x)\}(da)\,
 \Gamma(P,x)(d(y^0,y^1))\\
&\quad=\{1-e(P,x)\}\int_{\{0,1\}^2}
 \varphi(x,0,y^0)\,\Gamma(P,x)(d(y^0,y^1))\\
&\qquad\quad
 +e(P,x)\int_{\{0,1\}^2}
 \varphi(x,1,y^1)\,\Gamma(P,x)(d(y^0,y^1))\\
&\quad=\{1-e(P,x)\}
 \bigl[\{1-\mu_0(P,x)\}\varphi(x,0,0)
       +\mu_0(P,x)\varphi(x,0,1)\bigr]\\
&\qquad\quad
 +e(P,x)
 \bigl[\{1-\mu_1(P,x)\}\varphi(x,1,0)
       +\mu_1(P,x)\varphi(x,1,1)\bigr]\\
&\quad=\sum_{a=0}^1\sum_{y=0}^1p_{ay}(P,x)\varphi(x,a,y).
\end{aligned}
\]
Integrating over \(\lambda\), successive kernel integration and the observed-law cell disintegration give
\[
\int\varphi\circ\mathsf o_{\mathrm{obs}}\,
 d\widetilde P_{P,\Gamma}
=
\int_{[0,1]}\sum_{a=0}^1\sum_{y=0}^1
 p_{ay}(P,x)\varphi(x,a,y)\,\lambda(dx)
=
\int\varphi\,dP.
\]
Equality for all such tests proves
\[
(\mathsf o_{\mathrm{obs}})_\#\widetilde P_{P,\Gamma}=P.
\]
% lean: CausalSmith.Stat.LogoddsLowsmoothFrontier.causalExtension_observed_marginal

\item For a full-data record, the two treatment cases give
\[
Y=
\begin{cases}
Y^0,&A=0,\\
Y^1,&A=1,
\end{cases}
\qquad\text{and hence}\qquad
Y=(1-A)Y^0+AY^1.
\]
The successive-draw construction is exactly that of \cref{obj:def:causal-extension}. At every \(x\), its full conditional law is the product probability measure
\[
\mathcal K_{P,\Gamma}(x)
=\Gamma(P,x)\otimes\operatorname{Bernoulli}\{e(P,x)\}.
\]
The coordinate projections of a product probability measure are independent, so
\[
(Y^0,Y^1)\ \perp\ A\mid X=x
\qquad(x\in[0,1]).
\]
% lean: CausalSmith.Stat.LogoddsLowsmoothFrontier.observedMap_consistency
% lean: CausalSmith.Stat.LogoddsLowsmoothFrontier.fullConditionalLaw_independent

\item Finally, pushing the coordinate integrals through the prescribed Bernoulli margins yields
\[
\begin{aligned}
\int_{\{0,1\}^2}y^0\,\Gamma(P,x)(d(y^0,y^1))
&=0\{1-\mu_0(P,x)\}+1\mu_0(P,x)=\mu_0(P,x),\\
\int_{\{0,1\}^2}y^1\,\Gamma(P,x)(d(y^0,y^1))
&=0\{1-\mu_1(P,x)\}+1\mu_1(P,x)=\mu_1(P,x).
\end{aligned}
\]
Both means are interior. Substitution into the homogeneous log-odds identity established above gives, at every \(x\in[0,1]\),
\[
\operatorname{logit}
 \left(\int_{\{0,1\}^2}y^1\,\Gamma(P,x)(d(y^0,y^1))\right)
-
\operatorname{logit}
 \left(\int_{\{0,1\}^2}y^0\,\Gamma(P,x)(d(y^0,y^1))\right)
=\theta(P).
\]
% lean: CausalSmith.Stat.LogoddsLowsmoothFrontier.admissibleCoupling_means
% lean: CausalSmith.Stat.LogoddsLowsmoothFrontier.homogeneous_logit_difference
\end{enumerate}
\end{proof}
